\section{Positive Results}
\label{section_positive_results}

In this section we prove some positive results, including the facts that the LP for topologies with $n \ge 3$ nodes always has integer vertices that are not induced by STTs, and that the LP for a star graph has only integer vertices, implying optimality of the LP for stars of any size.

\subsection{Basic Properties}
\label{subsection_basic_properties_proven}

\begin{lemma}[Depths Bound]
\label{lemma_depth_bounds}
Let $(X,Z,D)$ be a feasible solution of the LP. Then for \emph{every} subset of coordinates $S \subseteq \{1,\ldots,n\}$ it holds that: $\sum_{i \in S}{D_i} \ge |S|-1$.
\end{lemma}
The claim is obvious for points induced by STTs, since the LP-depth of every node except for the root is at least one. For the general case, we argue as follows.

\begin{proof}
Fix $i$ and $j$ and consider their ancestry constraint:
$
1
\le X_{ij} + X_{ji} + \sum_{k \in (i \leftrightsquigarrow j)}{Z_{kij}}
\le X_{ij} + X_{ji} + \sum_{k \in (i \leftrightsquigarrow j)}{\min\{X_{ki},X_{kj}\}}
\le X_{ij} + X_{ji} + \frac{1}{2} \sum_{k \in (i \leftrightsquigarrow j)}{(X_{ki}+X_{kj})}
\le \frac{1}{2} (X_{ij} + X_{ji} + D_i + D_j)
$. Then: $\binom{|S|}{2} \le \sum_{\{i,j\} \subset S}{\frac{1}{2} (X_{ij} + X_{ji} + D_i + D_j)}$. Note that each $D_i$ appears exactly $|S|-1$ times on the right side, and that we have one more contribution of (at most) $D_i$ from the variables $X_{ji}$. Therefore: 
$\binom{|S|}{2} \le \frac{|S|}{2} \sum_{i \in S}{D_i} \Rightarrow |S|-1 \le \sum_{i \in S}{D_i}$.
\end{proof}

\begin{corollary}
\label{corollary_upper_bound_small_depth_counts}
The depths-vector of any feasible solution has at most $k$ entries with value at most $\frac{k}{k+1}$. Particularly, at most one coordinate smaller than $\frac{1}{2}$, at most two smaller than $\frac{2}{3}$, etc.
\end{corollary}

\begin{remark} 
Corollary~\ref{corollary_upper_bound_small_depth_counts} interprets Lemma~\ref{lemma_depth_bounds} as an upper-bound on the number of nodes with small depth $d < 1$ (it is not meaningful for $d \ge 1$). However, there are no lower-bound guarantees. Unlike STT induced vertices that are guaranteed to have one $0$ depth (the root's), a general vertex could have only ``large'' coordinates.

It may be interesting to understand whether the minimum depth (coordinate) of a vertex is bounded from above. We emphasize that this question is interesting only with respect to vertices rather than with respect to general feasible solutions, because it is easy to take a combinations of STTs to get a feasible point whose minimum depth is large. As a concrete example of vertices whose minimum depth is $1$, there are four such vertices up to automorphisms for the topology $U_{(8,4)}$ (Figure~\ref{figure_all_trees_up_to_n8}) which are $(2, 2, 4.5, 1.5, 1, 1, 2, 2)$, $(1.5, 2, 4.5, 1.5, 1, 1, 2.5, 2.5)$, $(2.5, 2.5, 4.5, 1.5, 1.5, 1.5, 1.5, 1)$, $(1.5, 2, 5, 1.5, 1, 1, 2.5, 2.5)$ and one vertex for the topology $U_{(8,13)}$ which is $(1.5, 2, 3.5, 2, 1, 1, 3.5, 1.5)$.
\end{remark}

Lemma~\ref{lemma_depth_bounds} is agnostic to weights. The following Lemma~\ref{lemma_node_depth_by_weight} takes weights into account.

\begin{lemma}[Depth by weights]
\label{lemma_node_depth_by_weight}
Let $i$ be a node whose frequency is $f_i$ where $\sum_{u \in U}{f_u} = 1$. Then $depth(i) \le \frac{1}{f_i}$ in every optimal STT. (Recall that the root's depth is $1$.)
\end{lemma}
\begin{proof}
Consider an STT $T_1$ where $i$ has depth $k > \frac{1}{f_i}$. Define $T_2$ by promoting $i$ to the root. Each node other than $i$ gains at most one new ancestor ($i$), so we lose a cost of at most $\sum_{j (\ne i) \in U}{f_j} = 1 - f_i$. However, we also save a cost of $(k-1) \cdot f_i$ for $i$. In total: $cost(T_1) - cost(T_2) \ge (k-1) f_i - (1- f_i) = k \cdot f_i - 1 > 0$. So $T_1$ is sub-optimal.
\end{proof}

\begin{corollary}
If there is a node $i$ with frequency $f_i > \frac{1}{2}$, it must be the root of every optimal STT. Therefore, we can focus our efforts on inputs where $\forall i: f_i \le \frac{1}{2}$. 
\end{corollary}



\subsection{More Analysis of Integer Vertices}
\label{subscetion_note_on_STT_as_vertices}
In this section we motivate focusing our interest on the lower envelope of the $D$-space polytope as we discussed in Section~\ref{subsection_general_counter_exaple} for the normals method.

\begin{theorem}
\label{theorem_integer_vertex_iff_stt}
Let $\mathbb{P}$ be the polytope of the LP corresponding to a topology $U$, and let $P \in \mathbb{P}$ be an integer point. If $P$ is induced by some STT $T$ then $P$ is a vertex on the lower envelope of $\mathbb{P}$. Moreover, if the projection of $P$ to $D$-space is a vertex on the lower envelope of the projection of $\mathbb{P}$, then $P$ is induced by some STT.
\end{theorem}

Note that Theorem~\ref{theorem_integer_vertex_iff_stt} is almost an ``if and only if'' statement, except that the first part of the claim considers the $(X,Z,D)$-polytope while the second part considers the projection to $D$-space.

\begin{proof}
We distinguish projections to $D$-space from the corresponding objects 
in the $(X,Z,D)$-polytope by a subscript $D$. Consider the second part of the claim, that $P_D$ is a vertex on the lower-envelope of $\mathbb{P}_D$, and assume by contradiction that $P$ is not induced by some STT. Then by being a vertex of $\mathbb{P}_D$, there is a direction $w$ for which $P$ is the unique optimum of the LP. Since $P_D$ is on the lower envelope of $\mathbb{P}_D$, there is such $w$ with non-negative coordinates. This $w$ is a natural objective function to Problem~\ref{problem_main_problem_stt}, and from Property~\ref{property_integer_domination} we deduce that $P$ is not the unique optimum in the direction $w$, a contradiction.

For the first part of the claim, define the weight of each node $i$ as $w_i = n^{-4 \cdot depth_T(i)}$ (recall that $T$ induces $P$). We next show that $P$ is the unique optimal solution for the LP in the direction $w$. This implies that $P$ is a vertex, and since all the weights are positive, that $P$ is on the lower envelope of $\mathbb{P}$.

To show that $P$ is uniquely optimal in the direction $w$, we prove that every optimal solution satisfies $\forall i: X_{ir}=0$ (this holds for $P$). Then $X_{ir}=0 \Rightarrow Z_{ijr} = 0 \; (\le X_{ir})$ and we deduce from the ancestry constraint on $r$ and $j$ that $\forall j: X_{rj} = 1$. The argument is concluded by recursively considering the connected components of the topology when $r$ is removed, and eventually after we fix all the variables, the set of optimal solutions turns out to be a single point, $P$.

Consider a feasible solution with $X$ coordinates such that there is some $i \ne r$ such that $X_{ir} > 0$, and let $\Delta \equiv \min \{ X_{ir} \mid X_{ir} > 0 \}$. We can infer values of  the other variables by setting: $D_i = \sum_{j}{X_{ji}}$ and $Z_{kij} = \min\{X_{ki},X_{kj}\}$.\footnote{There could be other ways to set $Z$ (smaller) and $D$ (larger), but our choice is feasible. Importantly, we cannot choose smaller values of $D$.} We define an alternative feasible solution $X'$ where $X'_{ij} = X_{ij} + n\Delta$ for $j \ne r$, $X'_{ir} = \max\{0,X_{ir} - \Delta\}$ and $Z'_{kij} = \min\{ X'_{ki},X'_{kj} \}$. If $j=r$ then $Z'_{kir} \ge Z_{kir} - \Delta$ ($Z'_{kir} = Z_{kir}$ if $X_{kr} = 0$), and if $r \ne i,j$ then $Z'_{kij} = Z_{kij} + n\Delta$.

By construction we maintain the non-negativity of $X'$ (and hence of $Z'$). Ancestry constraints are preserved because all the variables increase except for those of the form $X_{ir}$ and $Z_{kir}$ ($i \ne r$). $X_{ir}$ and $Z_{kir}$ only decrease if $X_{ir} > 0$, and can only violate the ancestry constraints over the pair $r$ and $i$. However, this constraint is not violated since: $X'_{ri} + X'_{ir} + \sum_{k \in (i \leftrightsquigarrow r)}{Z'_{kir}} \ge 
(X_{ri}+n\Delta) + (X_{ir}-\Delta) + \sum_{k \in (i \leftrightsquigarrow r)}{(Z_{kir} - \Delta)} \ge 1 + \Delta$. Finally, observe that $D'_r \le D_r - \Delta$ and for $i \ne r$: $D'_i \le D_i + n \cdot n\Delta$. Therefore the difference in values of the LP is:
$value(X',w) - value(X,w) = \sum_{i}{w_i \cdot (D'_i - D_i)} \le \sum_{i \ne r}{\frac{w_i}{w_r} \cdot w_r \cdot n^2 \Delta} - w_r \Delta < (n \cdot \frac{1}{n^4} \cdot n^2 - 1) \cdot w_r \Delta  < 0$. Therefore, every optimal point with respect to $w$ satisfies that $\forall i: X_{ir} = 0$.
\end{proof}


Theorem~\ref{theorem_integer_vertex_iff_stt} does not rule out an integer vertex $P$ of the LP such that $P_D$ is not on the lower envelope of $\mathbb{P}_D$. Indeed, Theorem~\ref{theorem_non_tree_vertex} shows that such vertices exist. Conveniently, their projection does not affect our normals method since we relax $\mathbb{P}_D$ to only deal with lower-envelope facets.

\begin{theorem}
\label{theorem_non_tree_vertex}
For any $n \ge 3$, and any tree topology $U$ with $n$ nodes, the LP defined by $U$ has integer vertices that are not induced by an STT (Definition~\ref{definition_xzd_of_stt}).
\end{theorem}

The proof shows two generic classes of such integer vertices. The first relies on non-transitivity in ancestral relationship, in opposition to an STT property not captured by the LP. The second defines an LCA implicitly but not explicitly.

\begin{proof}
Given a topology $U$, Let $\mathbb{P}'$ be the set of all feasible solutions such that every variable is $0$ or $1$, and in each ancestry constraint, exactly one variable is $1$. We argue that every point in $\mathbb{P}'$ is a vertex. Let $P \in \mathbb{P}'$.
$P$ is a vertex if and only if there is no non-zero direction $R$ such that $P \pm \epsilon R$ are both feasible. If $R$ is nonzero in some zero coordinate of $P$, one direction is infeasible due to negativity. Similarly, $R$ cannot be nonzero in  a coordinate where $P$ is $1$, because in order to preserve the ancestry feasibility to which this coordinate is the sole  contributor, $R$ must be non-zero on some other zero coordinate of $P$ with opposite sign. But since we already argued that $R$ is zero in all the zero coordinates of $P$, no such direction $R$ exists and we conclude that $P$ is a vertex.

Let $\mathbb{P} \subseteq \mathbb{P}'$ contain all the feasible solutions such that
$Z_{kij}=0$ for all $i$, $j$, and $k\in (i \leftrightsquigarrow j)$, and for every pair $i,j$ choose $X_{ij},X_{ji} \in \{0,1\}$ such that $X_{ij} + X_{ji} = 1$. We define two classes of points in $\mathbb{P}$ which do not correspond to STTs.

\begin{enumerate}
    \item We choose $P \in \mathbb{P}$ such that there is a triplet of indices $i,j,k$ (this is why the claim requires $n \ge 3$) for which $X_{ij} = X_{jk} = X_{ki} = 1$. No STT has such a cyclic relation, because if $i$ is an ancestor of $j$ which is an ancestor of $k$, then $i$ is also an ancestor of $k$.

    \item Start from a point $P' \in \mathbb{P}'$ induced by an STT $T$ that has at least one node that is an LCA ($T$ is not a path). Denote the LCA by $a$ and the two nodes that it separates by $b$ and $c$. Define $P$ to have the same values as $P'$ except that we set $Z_{abc} = 0$ and $X_{bc}=1,X_{cb}=0$. One can verify that $P$ is feasible.
    We say that $a$ remains an implicit ancestor of $b$ and $c$ since we still have that $X_{ab}=1$, $X_{ac}=1$, but 
    $Z_{abc} = 0$.
    If we replace every variable $Z_{kij}=1$ as we did for $Z_{abc}$, we  get a point in $\mathbb{P}$. Note that as a side-effect, we showed some non-STT integer vertices that belong to $\mathbb{P}'$ and not to $\mathbb{P}$ (some ancestors were made implicit but not all of them). \qedhere
\end{enumerate}
\end{proof}

\begin{remark}
\label{remark_most_integer_vertices_are_non_stt}
Theorem~\ref{theorem_non_tree_vertex} shows the existence of non-STT, even integer, vertices. In fact, a simple counting argument shows that they are the majority for large $n$.
Let $S$ be the set of all STT vertices, and $\mathbb{P}$ and $\mathbb{P}'$ as in the proof of Theorem \ref{theorem_non_tree_vertex}. We have that $|\mathbb{P}'| > |\mathbb{P}| = 2^{\binom{n}{2}}$, while $|S| \le n!$ since a permutation uniquely defines an STT by choosing a root and recursing on the suffix (multiple permutations may define the same STT). Finally, note that $S \subset \mathbb{P}'$.
\end{remark}


To conclude this section, we analyze all the vertices of the polytopes for topologies with $n=2$ and $n=3$ nodes, each of these sizes has a single, path, topology. In both cases, there are only integer vertices. When $n=2$, there are no $Z$ variables, the LP is defined as $X_{12}+X_{21} \ge 1$ for $X_{12},X_{21} \ge 0$. Overall there are two vertices, one per STT. When $n=3$ the LP becomes more interesting. The following inequalities define the polytope, and Table~\ref{table_vertices_n3} lists all $9$ vertices: $5$ of them correspond to STTs, and $4$ more are constructed as detailed in the proof of Theorem~\ref{theorem_non_tree_vertex}, two of each class.
\begin{equation*}
    X_{12}+X_{21} \ge 1, \ \ \ 
    X_{23}+X_{32} \ge 1,  \ \ \ 
    X_{13}+X_{31}+Z_{213} \ge 1,  \ \ \ 
    Z_{213} \le X_{21},  \ \ \ 
    Z_{213} \le X_{23}  \ \ \ 
\end{equation*}

\begin{table}[!ht]%[!ht]
    \scriptsize
    \begin{center}
    \begin{tabular}{|c|c|c|c|c|c|c|c|c|l|}
    \hline
    $X_{12}$ & $X_{21}$ & $X_{23}$ & $X_{32}$ & $X_{13}$ & $X_{31}$ & $Z_{213}$ & Comments \\
    \hline
    \hline
    0 & 1 & 0 & 1 & 0 & 1 & 0 & STT: root $3$, leaf $1$ \\
    \hline
    0 & 1 & 0 & 1 & 1 & 0 & 0 & non-STT: cyclic ancestry  \\
    \hline
    0 & 1 & 1 & 0 & 0 & 0 & 1 & STT: root $2$ \\
    \hline
    0 & 1 & 1 & 0 & 0 & 1 & 0 & non-STT: LCA abuse  \\
    \hline
    0 & 1 & 1 & 0 & 1 & 0 & 0 & non-STT: LCA abuse  \\
    \hline
    1 & 0 & 0 & 1 & 0 & 1 & 0 & STT: root $3$, leaf $2$ \\
    \hline
    1 & 0 & 0 & 1 & 1 & 0 & 0 & STT: root $1$, leaf $2$  \\
    \hline
    1 & 0 & 1 & 0 & 1 & 0 & 0 & STT: root $1$, leaf $3$  \\
    \hline
    1 & 0 & 1 & 0 & 0 & 1 & 0 & non-STT: cyclic ancestry \\
    \hline
    \hline
    \end{tabular}
    \end{center}
    \caption{\footnotesize{All the vertices of the LP of the topology over $n=3$ nodes. The vertices with comment of ``non-STT: LCA abuse'' become non-vertices when we eliminate $Z_{213}$ as detailed in Section~\ref{subsection_with_and_without_z_fourier_motzkin}.}}
    \label{table_vertices_n3}
\end{table}



\subsection{Star Topologies Have Integer Vertices}
\label{subsection_star_is_integer}

In this subsection we prove analytically that if the tree topology is a star, then all the vertices of the LP are integer. Section~\ref{subscetion_note_on_STT_as_vertices} explicitly shows this claim for $n=2,3$. Consequently, as a result of Property~\ref{property_integer_domination}, we get a polynomial-time algorithm to find the optimal static STT (for stars). 

\begin{theorem}
\label{theorem_star_has_integer_vertices}
The LP polytope for a star topology has only integer vertices.
\end{theorem}

\begin{corollary}
An optimal static STT for any star topology is computable in polynomial time.
\end{corollary}

(The corollary can be proven directly, using a simple enumeration: since the star is very symmetric, there are at most $n$ effective STTs to consider: we always query leaves in decreasing order of weight, and the $n$ STTs vary according to the depth at which we query the central node.)



\begin{proof}[Proof of Theorem~\ref{theorem_star_has_integer_vertices}]
The LP for a star over $n$ nodes is given in Equation~(\ref{equation_starn_original}), where we denote the center by $1$ (recall that $Z_{1ij}$ and $Z_{1ji}$ are the same).
\begin{equation}
\label{equation_starn_original}
    \forall i>1: X_{1i}+X_{i1} \ge 1 ; \\  \ \ \ \ \ 
    \forall i,j \ne 1: X_{ij} + X_{ji} + Z_{1ij} \ge 1 ; \\  \ \ \ \ \ 
    \forall i,j \ne 1: Z_{1ij} \le X_{1i} ; \\  \ \ \ \ \ 
    X,Z \ge 0 \\
\end{equation}

We show that each vertex corresponds to an assignment satisfying the following three conditions:
\begin{enumerate}
    \item \label{assign1} For every $i,j \ne 1$, $Z_{1ij} \in \{0,1\}$.
    \item \label{assign2} For each $i$: Let $z_i \equiv \max_{j \ne 1,i} \{Z_{1ij}\}$. We set $X_{1i} \in \{z_i,1\}$ and  $X_{i1} = 1 - X_{1i}$.
    \item \label{assign3} For each $i,j \ne 1$: Choose $X_{ji} \in \{0,1-Z_{1ij}\}$ and set $X_{ij}=1-Z_{1ij}-X_{ji}$. Explicitly: if $Z_{1ij}=1$ then $X_{ij}=X_{ji}=0$, otherwise ($Z_{1ij} = 0$) $X_{ij} \in \{0,1\}$ and $X_{ji} = 1-X_{ij}$.
\end{enumerate}
One can verify that each such assignment defines a feasible integer point. We proceed to prove that every vertex corresponds to an assignment that satisfies these conditions. We do this by showing that if any of the three conditions does not hold at a feasible point $P$, then there is a direction $R$ such that $P'_{\pm} = P \pm \epsilon R$ are both feasible for some $\epsilon > 0$, thus $P = \frac{1}{2} (P'_{+} + P'_{-})$ is not a vertex. The coordinates of $R$ are either $0$ or $\pm 1$, and we say that a coordinate of $R$ 
that equals $1$ is \emph{increased}
and a coordinate of $R$ that equals $-1$ is \emph{decreased}.
 Once we prove that some condition holds, we assume it to hold when we consider other conditions in later arguments.

First note that every variable is bounded in $[0,1]$. Non-negativity is by definition of the LP. If there are variables larger than $1$, let $R$ decrease all the variables that are larger than $1$ (indeed $P \pm \epsilon R$ are both feasible).

To prove Condition~\ref{assign3}, for fixed $i,j \ne 1$ note that these variables only participate in the constraint $X_{ij}+X_{ji} \ge 1 - Z_{1ij}$. If $X_{ij}+X_{ji} > 1 - Z_{1ij}$, since $Z_{1ij} \le 1$ one of $X_{ij}$ and $X_{ji}$ is positive, and we can choose $R$ to decrease it. Thus $X_{ij} + X_{ji} = 1 - Z_{1ij}$. If $X_{ij},X_{ji} \in (0,1-Z_{1ij})$ (note that if $Z_{1ij}=1$ this range is empty) then both are positive, and $R$ increases $X_{ij}$ and decreases $X_{ji}$. We conclude that $X_{ij}$ is either $0$ or $1-Z_{1ij}$ and $X_{ji} = 1-Z_{1ij}-X_{ij}$. This proves Condition~\ref{assign3}.

To prove Condition~\ref{assign2}, consider the case $X_{i1} + X_{1i} > 1$ for some $i$. Since both variables are at most $1$, both are positive, and $X_{i1}$ only participates in this single inequality, so take $R$ which increases $X_{i1}$. Now that we have $X_{i1} + X_{1i} = 1$, if $z_i < X_{1i} < 1$ then $X_{i1} > 0$, so take $R$ which decreases $X_{1i}$ and increases $X_{i1}$. (In both cases, the strict inequality ensures that $P \pm \epsilon R$ are both feasible, so we must have equalities as argued.)

To prove Condition~\ref{assign1}, let $0<a<1$ be the largest non-integer value of some $Z$ variable, if such exists. For each $Z_{1ij}=a$ we have $(X_{ij}+X_{ji}) \ge (1-a) > 0$, so let $A_{ij}$ denote a positive variable among $X_{ij}$ and $X_{ji}$ (if both are, pick one arbitrarily). Then $R$ increases all the $X$ and $Z$ variables whose value is in $[a,1)$, decreases the variables $A_{ij}$ that correspond to an increased $Z_{1ij}$, and if $X_{1i}$ was increased (because $a \le X_{1i} < 1$), recall that since Condition~\ref{assign2} holds, $X_{1i} + X_{i1} = 1$, so we also decrease $X_{i1}$ to maintain feasibility in the direction $-R$. $P \pm \epsilon R$ are both feasible, so if $P$ is a vertex, every $Z_{1ij} \in \{0,1\}$ as required by Condition~\ref{assign1}.
\end{proof}



\subsection{Extending Topologies}
\label{subsection_graph_extensions}

In this section we prove that if the LP for a topology has non-integer vertices, the LP of topologies that extend it in some ways also have non-integer vertices. This explains formally, for example, why the three topologies with $8$ nodes that are colored in red in Figure~\ref{figure_all_trees_up_to_n8} must have non-integer vertices once we know that the red topology with $7$ nodes has non-integer vertices. We then discuss another kind of extension, when we combine topologies through a new common node.

\begin{definition}
\label{definition_extended_graph}
We define a single extension of a graph $G$ to $G'$ and denote it $G \to G'$ as the action of adding a new node to $G$ and connecting it to some existing node as a new leaf, or between two adjacent nodes, i.e., subdividing an existing edge into two new edges.
We say that $G$ extends $H$, or that $H$ is a shrinking of $G$, if there is a sequence of extensions of length $|G|-|H|$ such that $H \to H' \to \ldots \to G$.
\end{definition}

\begin{theorem}
\label{theorem_subgraph_with_frac_vertex_implies_frac_vertex_extension}
If the LP for a topology $U$ has non-integer vertices, so does the LP for any extension of $U$. The claim holds both for the original space of $(X,Z,D)$ variables, and also when only considering the depths $D$.
\end{theorem}


\begin{proof}
It is sufficient to prove the claim for a single extension from $U$ to $U'$. Then if $U^*$ is an extension of $U$, by applying the claim over a single extension $|U^*|-|U|$ times, we may conclude that $U^*$ too has non-integer vertices.


Throughout the proof we  use prime to denote terms that correspond to $U'$. For example, we denote the set of variables in the LP of $U$ by $(X,Z,D)$, and  $(X',Z',D')$ denote the variables in the LP of $U'$. Since $U'$ is a single extension of $U$, $(X',Z',D')$ is a union of $(X,Z,D)$ with additional new variables that correspond to the new extending node, which we denote by $a$. 

Let $h$ be a direction for which a fractional vertex $P$ is uniquely optimal for the $LP$ of $U$ (there is such direction for any vertex). Let $h'$ extend $h$ with $0$'s for every new coordinate, and let $Q'$ be an optimal vertex for the $LP$ of $U'$ (it may not be uniquely optimal, but this is fine). We will show that we can extend $P$ to a point $P'$ feasible for the LP of $U'$ such that $P$ is the projection of $P'$ to the old variables. Denote by $Q$ the projection of $Q'$ to the old variables. Observe that: $Q \cdot h = Q' \cdot h' \le P' \cdot h' = P \cdot h$ where both equalities are because $h'$ has $0$'s in coordinates of new variables, and the inequality follows since $Q'$ is a minimizer in the direction $h'$. Because $P$ is uniquely optimal for $h$, then $Q = P$, and we conclude that $Q'$ is the vertex claimed by the theorem: $Q'$ projects to $P$, hence if $P$ is non-integer so is $Q'$. Moreover, if $P$ has non-integer $D$ variables, so does $Q'$.

We explain how to extend $P$ to $P'$ that is a feasible point for the LP of $U'$, not necessarily a vertex, by setting the values of the extra variables. Loosely speaking, we set the values so that $a$ is a ``descendant'' of its neighbor(s). The values of $X,Z,D$ are the same in $P'$ and $P$, and we only need to define the values of $X'\setminus X$, $Z' \setminus Z$ and $D_a (= D'\setminus D)$. Because $U'$ is a single extension, we have two cases to consider. The first and simpler case is when the added node $a$ is a leaf, denote its neighbor by $b \in U$, and $U_b \equiv U \setminus \{b\}$. Then the extension is as follows:
\begin{enumerate}
    \item $X' \setminus X$: Set $X'_{ba}=1$, $\forall i \in U_b: X'_{ia} = X_{ib}$ and $\forall i \in U: X'_{ai} = 0$ (child of $b$, ancestor of none).
    \item $Z' \setminus Z$: Set $Z'_{kai} = Z_{kbi}$ for $i,k \in U_b$ and $Z'_{bai} \equiv \min\{ X'_{ba}, X'_{bi} \} = X_{bi}$. Since $a$ is a leaf in the topology there are no variables of the form $Z'_{aij}$.
    \item Set $D'_a = \sum_{i \in U}{X'_{ia}}$.
\end{enumerate}

To verify feasibility, note that any constraint involving only old variables is satisfied by construction because we copied the values. So we only need to check constraints involving variables associated with $a$. By construction $D'_a \ge \sum_{i \in U}{X'_{ia}}$ (in fact, equal), and since $\forall i \ne a: X'_{ai}=0$ we have indeed $D'_i = D_i$. Next, $Z'_{iaj} \le \min\{ X'_{ia},X'_{ij} \}$ is explicitly stated if $i=b$, and for $i \ne b$: $Z'_{iaj} = Z_{ibj} \le \min\{ X_{ib},X_{ij} \} = \min\{ X'_{ia},X'_{ij} \}$. Finally,
consider ancestry constraints for $a$ and some $i \in U$. If $i=b$: Then $X'_{ba} = 1$ and we are done. Otherwise, $X'_{ai} + X'_{ia} + \sum_{k \in (a \leftrightsquigarrow i)}{Z'_{kai}} = 0 + X_{ib} + (X_{bi} + \sum_{k \in (b \leftrightsquigarrow i)}{Z_{kbi}}) \ge 1$ where the last inequality is known to hold, in $U$, with respect to the pair $i,b$.

In the second case node $a$ subdivides the edge $(b,c)$ between two nodes $b,c \in U$. Define $U'_x$ for $x \in \{b,c\}$ to be the connected component of $U'$ that contains $x$ when we delete $a$. We set the $X$ and $Z$ variables that involve $a$ and nodes in $U'_x$ as if $a$ is a leaf that was added to $U'_x$, and set $D'_a = \sum_{i \in U}{X'_{ia}}$. It only remains to define values for the variables  $Z'_{aij}$ for $i \in U'_b$ and $j \in U'_c$ (a case we did not have earlier), and we trivially set them to $0$.

To verify feasibility, note that setting $Z'_{aij} = 0$ for $i \in U'_b$ and $j \in U'_c$ is fine because $Z'_{aij} = 0 \le X'_{ai} = X'_{aj}  = 0$. Then the only non-trivial constraints that require verification are ancestry constraints. However, for pairs $i,j \in U$ we have them satisfied by construction (copied values), and for pairs of $a$ with $i \in U'_x$ we know them to be satisfied because we verified this relation in the previous case of $a$ being an actual leaf (here $U'_y$ for $y \ne x$ has no effect on the constraints).
\end{proof}

\begin{remark}
\label{remark_subgraph_notes}
Note that the converse is of course not true: if an extension has non-integer vertices, it does not imply that the smaller topology has non-integer vertices. Indeed, the path topology with $3$ vertices only has integer vertices as shown in Table~\ref{table_vertices_n3}, but is extended by the long-star which has non-integer vertices (Theorem~\ref{theorem_non_integer_optimum_long_star}). The proof works since we can add zero coefficients when extending $h$ to $h'$.
\end{remark}


Next, we show that we can ``mix'' vertices of smaller topologies, in larger topologies.

\begin{theorem}
\label{theorem_combine_trees}
Let $b \ge 2$ be an integer, and let $H = \{ U^1,\ldots,U^b \} $ be a collection of topologies. Let $U$ be a new (tree) topology 
obtained by taking the union over $H$ and adding an additional new node $r$, connected to one node in each $U^i$.
 Then for every set of vertices $P^1,\ldots,P^b$ such that $P^i$ is a vertex of the LP for $U^i$, there is a vertex $P$ whose projection to variables that are only related to $U^i$ is $P^i$. In loose terms, $P$ is sort of a Cartesian product.
\end{theorem}

\begin{proof}
To define $P$, we intuitively set $r$ as the ``ancestor'' of every other node, which formally translate to the following values: $\forall u \in U' \setminus \{r\}: X_{ru} = 1, X_{ur} = 0$, and $\forall u \in U^i, v \in U^j, i \ne j: X_{uv}=X_{vu}=0$. Also, $D_r = \sum_{i \in U' \setminus\{r\}}{X_{ir}} = 0$. In addition, the LCA relations are as follows: $\forall u \in U' \setminus \{r\}: \forall k \in (u \leftrightsquigarrow r): Z_{kur}  = 0$ and $\forall u \in U^i, v \in U^j, i \ne j: \forall k \in (u \leftrightsquigarrow v): Z_{ruv}  = \delta_{rk}$ (i.e., $1$ iff $k=r$). So far we set  all the new variables, related to $r$ or to pair of nodes from different $U^i$s. Within each $U^i$, we copy the values of the vertex $P^i$. This results in a feasible point $P$ that is a ``Cartesian product'' of $P^1,\ldots,P^b$. Feasibility is simple to verify because $r$ behaves like a root, and within each subtree we have feasibility by definition. It remains to prove that $P$ is a vertex.

Assume by contradiction that $P$ is not a vertex, then there is a direction vector $R$ such that $P \pm \epsilon R$ is feasible. However, observe that the values that we defined explicitly are all extremal, either $0$ or $1$. Because we cannot decrease below $0$, the coordinates of variables which we set to $0$ in $P$ must be $0$ also in $R$. As for the ones that are $1$, we claim that they cannot be decreased. Decreasing $X_{ru}$ must be accompanied with increasing one of $X_{ur}$ or some $Z_{kur}$ where $k \in (u \leftrightsquigarrow r)$, but these are zeros and as we argued must be $0$ also in $R$. Similarly, decreasing some $Z_{ruv} = 1$ necessitates increasing some among $X_{uv}$, $X_{vu}$, or $Z_{kuv}$ for $k \ne r$, but again these are
$0$ in $P$ and must be $0$ in $R$. Finally, the rest of the coordinates of $R$ are partitioned to $b$ groups by the topology $U^i$ to which their corresponding vertices belong. Since $P$ restricted to coordinates of topology $j$ is a vertex in $U^j$ we conclude that $R = \vec{0}$ and that $P$ is a vertex. 
\end{proof}

\begin{remark}
\label{remark_combine_trees}
Note that Theorem~\ref{theorem_combine_trees} does not argue that every vertex of $U$ is such a ``Cartesian product'', and it is not true in general. As a concrete example, consider $U^1=U^2=U^3$ to be trees with $2$ nodes each. They only have integer vertices so every combination is integer. However, the long-star which combines them via its center node, has non-integer vertices (Theorem~\ref{theorem_non_integer_optimum_long_star}).
\end{remark}

%%% The following example has been commented out for the sake of space, even though it gives a slightly more concrete intuition. The example shows that nodes with both thirds and halves coordinates exist.
% \begin{example}
% \label{example_mixed_halves_and_thirds}
% We show in Section~\ref{subsection_with_and_without_z_fourier_motzkin} that the path topology with $n=5$ nodes has integer, half-integer, and third-integer vertices in $(X,Z,D)$-space. Then the path with $n=11$ nodes also has vertices that are mixed with halves and thirds. Indeed, when rooting the sixth node (by choosing its frequency to be very large compared to the rest), we get two independent sub-STTs and we can choose the rest of the frequencies such that one sub-STT contributes half-integers and the other contributes third-integers.
% \end{example}

To conclude this section, we close a small debt from Section~\ref{subsection_general_counter_exaple} where we mentioned that primary directions may have additional similarities that are not automorphism symmetries.

\begin{remark}
\label{remark_pseudo_symmetry}
Consider two different shrinkings of $U$ (see Definition~\ref{definition_extended_graph}), denoted by $U_1$ and $U_2$, such that they are isomorphic and each has a different  subset of the nodes of the original $U$ ($U_1 \ne U_2$). Denote the isomorphism between $U_1$ and $U_2$ by $\pi$. By Theorem~\ref{theorem_subgraph_with_frac_vertex_implies_frac_vertex_extension}, for every vertex $P_i$ of $U_i$ ($i=1,2$) we have a vertex $P'_i$ of $U$ such that the coordinates of the variables that correspond to $U_i$
in $P_i$ and $P'_i$ are the same. 
Furthermore, the proof shows how to extend a direction $h_i$ for which $P_i$ is uniquely optimal, to a direction $h'_i$ for which $P'_i$ is optimal. The extension adds $0$'s in coordinates of $h'_i$ that do not correspond to coordinates of $h_i$. So if $P_1$ is a vertex of the LP of $U_1$ and $h_1$ is a direction in which $P_1$ it is the unique optimum, then $P_2 = \pi(P_1)$ is a vertex of $U_2$ and $h_2 = \pi(h_1)$ is a direction in which $P_2$ is uniquely optimal.
It follows that $P_1'$ and $P_2'$ are both vertices of $U$, $P_1'$ is optimal in the direction $h'_1$, and $P_2'$ is optimal in the direction $h'_2$. 
Note that
$h'_1$ and $h'_2$ are directions with the same coordinate-values, permuted. These values are the same as of $h_1$, and  $h_2$, padded with zeros. In many cases, this permutation does not correspond to an automorphism of $U$. As an example, name the edges of the topology $U_{(8,4)}$ as $\{(1,2),(2,3),(3,4),(4,5),(5,6),(3,7),(7,8)\}$. We can shrink it to $U_{(7,3)}$ by deleting one of the nodes in $\{4,5,6\}$, and as a result we get multiple similar but not automorphically symmetric primary directions. A concrete example of such directions is: $(3,2,0,0,2,3,3,10)$, $(3,2,0,2,0,3,3,10)$ and $(3,2,0,2,3,0,3,10)$ (all extend $(3,2,0,2,3,3,10)$ from Theorem~\ref{theorem_non_integer_optimum_long_star}). Such directions may find vertices that are not automorphically symmetric, and thus lead to a different behavior under rounding (e.g. with Definition~\ref{definition_rounding_scheme}).
\end{remark}